% ============================================================
% CONFIGURAÇÃO DO ARTIGO
% Disciplina XRSC10
%
% Professor: Pioli
% ============================================================


% ============================================================
% INFORMAÇÕES GERAIS
% ============================================================

\titulo{Análise Comparativa de FedAvg e FedProx sob Diferentes Tipos e Níveis de Heterogeneidade de Dados: Acurácia, Convergência e Custo Computacional}

\linkLatex{https://github.com/el-yawd/seminar-iv}


\ano{2026}

\mes{Dezembro}


% ============================================================
% AUTOR 1
% ============================================================

\autorNomeA{Diego}

\autorSobrenomeA{Reis}

\emailAutorA{d2022014641@unifei.edu.br}


% ============================================================
% AUTOR 2
% ============================================================

\autorNomeB{Augusto}

\autorSobrenomeB{Lázaro}

\emailAutorB{d2018006160@unifei.edu.br}


% ============================================================
% DISCIPLINA
% ============================================================

\disciplina{XRSC10}


% ============================================================
% PROFESSOR
% ============================================================

\professor{Pioli}


% ============================================================
% PALAVRAS-CHAVE
% ============================================================

\palavraChaveA{Aprendizado Federado}

\palavraChaveB{Computação de Borda}

\palavraChaveC{Dados Não-IID}

\palavraChaveD{FedAvg}

\palavraChaveE{FedProx}


% ============================================================
% COMANDOS E PACOTES DO TRABALHO
% Usados pelas seções 2 a 7 (tabelas, figuras, códigos e
% números); a classe PioliEdgeComp.cls não é alterada.
% ============================================================

% Link do repositório do projeto, exibido no resumo
\newcommand{\linkProjeto}{\url{https://github.com/el-yawd/seminar-iv}}

% Marca trechos que dependem de resultados ainda não disponíveis.
% Remover antes da entrega da versão final.
\newcommand{\pendente}[1]{\textcolor{red}{[#1]}}

\usepackage{booktabs}

\usepackage{siunitx}
\sisetup{output-decimal-marker={,}, group-separator={.}}

\usepackage{tikz}
\usetikzlibrary{positioning, arrows.meta, shapes.geometric}

\lstset{
    basicstyle=\scriptsize\ttfamily,
    numberstyle=\scriptsize\ttfamily,
    columns=fullflexible,
    keepspaces=true,
    literate={á}{{\'a}}1 {ã}{{\~a}}1 {ç}{{\c c}}1 {é}{{\'e}}1 {ê}{{\^e}}1
             {í}{{\'i}}1 {ó}{{\'o}}1 {õ}{{\~o}}1 {ú}{{\'u}}1
}
